% from ega2-1-fr.tex
\oldpage[II]{67}

On notera que cette condition est toujours vérifiée si $\psi$ est
(\textbf{TN})-surjectif.

De même, la question de savoir si $r_{\mathcal{L},\psi}$ est \emph{dominant}
est locale sur $Y$, et l'on a~:

\begin{corollary}[3.7.5]
\label{II.3.7.5-fr}
Sous les hypothèses de (\hyperref[II.3.7.3-fr]{3.7.3}), pour que
$r_{\mathcal{L},\psi}$ soit dominant, il faut et il suffit que, pour tout
entier $n>0$, toute section $s\in S_n$ telle que
$\psi^\flat(s)\in\Gamma(X,\mathcal{L}^{\otimes n})$ soit localement
nilpotente, soit elle-même nilpotente.
\end{corollary}

\begin{proof}
On doit en effet exprimer que
$r_{\mathcal{L},\psi}^{-1}(D_+(s))$ n'est vide que si $D_+(s)$ est vide,
et le corollaire résulte donc de
(\hyperref[II.3.7.3.1-fr]{3.7.3.1}) et de
(\hyperref[II.2.3.7-fr]{2.3.7}).
\end{proof}

\begin{proposition}[3.7.6]
\label{II.3.7.6-fr}
Soient $q:X\to Y$ un morphisme, $\mathcal{L}$ un $\mathcal{O}_X$-Module
inversible, $\mathcal{S}$, $\mathcal{S}'$ deux $\mathcal{O}_Y$-Algèbres
graduées quasi-cohérentes,
$u:\mathcal{S}'\to\mathcal{S}$ un homomorphisme d'Algèbres graduées,
$\psi:q^*(\mathcal{S})\to\bigoplus_{n\geq0}\mathcal{L}^{\otimes n}$ un
homomorphisme d'Algèbres graduées,
$\psi'=\psi\circ q^*(u)$ l'homomorphisme composé. Si
$r_{\mathcal{L},\psi'}$ est partout défini, il en est de même de
$r_{\mathcal{L},\psi}$~; si $u$ est (\textbf{TN})-surjectif et si
$r_{\mathcal{L},\psi'}$ est dominant, il en est de même de
$r_{\mathcal{L},\psi}$~; inversement, si $u$ est (\textbf{TN})-injectif et si
$r_{\mathcal{L},\psi}$ est dominant, $r_{\mathcal{L},\psi'}$ est dominant.
\end{proposition}

\begin{proof}
On a en effet $G(\psi')\subset G(\psi)$
(\hyperref[II.2.8.4-fr]{2.8.4}), d'où la première assertion~; si $u$ est
(\textbf{TN})-surjectif,
$\operatorname{Proj}(u):\operatorname{Proj}(\mathcal{S})\to
\operatorname{Proj}(\mathcal{S}')$ est partout défini et est une immersion
fermée~; comme $r_{\mathcal{L},\psi'}$ est composé de
$\operatorname{Proj}(u)$ et de la restriction de $r_{\mathcal{L},\psi}$ à
$G(\psi')$, on en conclut que si $r_{\mathcal{L},\psi'}$ est dominant, il en
est de même de $r_{\mathcal{L},\psi}$. Enfin, si $u$ est
(\textbf{TN})-injectif, on sait que $\operatorname{Proj}(u)$ est un morphisme
dominant de $G(u)$ dans $\operatorname{Proj}(\mathcal{S}')$
(\hyperref[II.2.8.3-fr]{2.8.3})~; comme $G(\psi')$ est l'image réciproque de
$G(u)$ par $r_{\mathcal{L},\psi}$, on voit que si $r_{\mathcal{L},\psi}$ est
dominant, il en est de même de $r_{\mathcal{L},\psi'}$.
\end{proof}

\begin{proposition}[3.7.7]
\label{II.3.7.7-fr}
Soient $Y$ un préschéma quasi-compact, $q:X\to Y$ un morphisme quasi-compact,
$\mathcal{L}$ un $\mathcal{O}_X$-Module inversible, $\mathcal{S}$ une
$\mathcal{O}_Y$-Algèbre graduée quasi-cohérente, limite inductive filtrante
d'un système inductif $(\mathcal{S}^\lambda)$ de $\mathcal{O}_Y$-Algèbres
quasi-cohérentes. Soient
$\varphi_\lambda:\mathcal{S}^\lambda\to\mathcal{S}$ l'homomorphisme
canonique,
$\psi:q^*(\mathcal{S})\to\bigoplus_{n\geq0}\mathcal{L}^{\otimes n}$ un
homomorphisme d'Algèbres graduées, et posons
$\psi_\lambda=\psi\circ q^*(\varphi_\lambda)$. Pour que
$r_{\mathcal{L},\psi}$ soit partout défini, il faut et il suffit qu'il existe
un $\lambda$ tel que $r_{\mathcal{L},\psi_\lambda}$ soit partout défini~;
$r_{\mathcal{L},\psi_\mu}$ est alors partout défini pour $\mu\geq\lambda$.
\end{proposition}

\begin{proof}
La condition est suffisante en vertu de
(\hyperref[II.3.7.6-fr]{3.7.6}). Inversement, supposons
$r_{\mathcal{L},\psi}$ partout défini~; on peut se ramener au cas où $Y$ est
affine, car si pour tout ouvert affine $U\subset Y$, il existe $\lambda(U)$
tel que la restriction de $r_{\mathcal{L},\psi_{\lambda(U)}}$ à $q^{-1}(U)$
soit partout définie, il suffira ($Y$ étant quasi-compact) de recouvrir $Y$
par un nombre fini d'ouverts affines $U_i$, et de prendre
$\lambda\geq\lambda(U_i)$ pour tous les indices $i$, en vertu de
(\hyperref[II.3.7.6-fr]{3.7.6}). Si $Y$ est affine, l'hypothèse entraîne que
pour tout $x\in X$, il y a une section $s^{(x)}$ d'un $S_n$ telle que, si on
pose $t^{(x)}=\psi^\flat(s^{(x)})$, on ait $t^{(x)}(x)\neq0$
($t^{(x)}$ étant considérée comme section de $\mathcal{L}^{\otimes n}$
au-dessus de $X$), ce qui entraîne $t^{(x)}(z)\neq0$ pour tout $z$ dans un
voisinage $V(x)$ de $x$. Couvrons $X$ par un nombre fini de $V(x_i)$ et
soient $s^{(i)}$ les sections correspondantes de $S$~; il existe alors un
$\lambda$ tel que les $s^{(i)}$ soient toutes de la forme
$\varphi_\lambda(s_\lambda^{(i)})$, avec $s_\lambda^{(i)}\in S^\lambda$ pour
tous les $i$~; il résulte donc de (\hyperref[II.3.7.4-fr]{3.7.4}) que
$r_{\mathcal{L},\psi_\lambda}$ est partout défini. La dernière assertion est
conséquence triviale de (\hyperref[II.3.7.6-fr]{3.7.6}).
\end{proof}

\begin{corollary}[3.7.8]
\label{II.3.7.8-fr}
Sous les hypothèses de (\hyperref[II.3.7.7-fr]{3.7.7}), si les
$r_{\mathcal{L},\psi_\lambda}$ sont dominants, il en est de même de
$r_{\mathcal{L},\psi}$~; la réciproque est vraie lorsque les
$\varphi_\lambda$ sont injectifs.
\end{corollary}

